\part{BaseCore Boundary and Expansion Interface}
\section{BaseCore Ownership and Downstream Boundary}

\begin{definition}[BaseCore Envelope]
The BaseCore envelope is the collection of theorem-level structures owned directly by this source tree:
\begin{enumerate}[leftmargin=*]
\item foundational Hilbert-space projection dynamics;
\item typed computable witnesses and state spectral-gap bounds;
\item parameterwise non-collapse under explicit anti-constant assumptions;
\item inverse-limit identity under non-finite-determination assumptions;
\item conditional rigidity under admissible perturbation models;
\item conditional non-simulability of CCR targets by finite simulator classes;
\item operational-criterion grammar, certification, falsation, and claim-boundary ledgers.
\end{enumerate}
\end{definition}

\begin{proposition}[Boundary Discipline]
BaseCore may define structures that downstream layers use, but it does not claim downstream closure. In particular, BaseCore does not by itself own runtime thresholds, bridge admissibility, human-comparator programs, phenomenality burdens, or Papers 7--9.
\end{proposition}

\begin{proof}
These excluded layers require additional assumptions, protocols, or empirical burdens not stated in BaseCore. Therefore they cannot be claimed as BaseCore theorems without violating the source-level claim boundary.
\end{proof}

\begin{center}
\renewcommand{\arraystretch}{1.12}
\begin{tabularx}{\textwidth}{P{0.31\textwidth} P{0.31\textwidth} Y}
\toprule
\textbf{Owned by BaseCore} & \textbf{Expanded but not owned here} & \textbf{Kept downstream} \\
\midrule
Projection dynamics, fixed points, compact attractors, anti-constant H5 burden, typed computable witness, inverse-limit identity under NFD, conditional rigidity, conditional non-simulability, operational-criterion grammar, falsation ledger &
Non-convex class variants, $\mathscr{C}^R$ style extensions, $\Omega$ dynamics, threshold numerics, mutable runtime class ledgers, implementation-specific status surfaces &
Paper 7 operational life / subjecthood, Paper 8 indexed subjectivity, Paper 9 phenomenal bridge organization, Stage 2 / Stage 3 comparative programs, human comparator families, phenomenality and external adjudication \\
\bottomrule
\end{tabularx}
\end{center}

\begin{remark}[Why the earlier material was removed from BaseCore ownership]
Previous versions mixed theorem-tight base results with runtime-like numerics, non-convex extension classes, and bridge-facing vocabulary. That blurred the mathematical base layer. BaseCore now keeps only the theorem ownership that survives the stricter release gate.
\end{remark}

\section{Admissible Structural Extensions}

\begin{definition}[BaseCore export object]
A BaseCore export object is a tuple
\[
  \mathfrak{E}=(\mathcal{O},\mathcal{H},\mathcal{R},\mathcal{B})
\]
where $\mathcal{O}$ is a typed mathematical object defined in BaseCore, $\mathcal{H}$ is the finite list of hypotheses under which it is used, $\mathcal{R}$ is the list of results actually proved from those hypotheses, and $\mathcal{B}$ is a boundary list of claims not entailed by $\mathcal{R}$.
\end{definition}

\begin{definition}[Downstream admissible extension]
A downstream layer is an admissible extension of a BaseCore export object $\mathfrak{E}$ only if it supplies:
\begin{enumerate}[leftmargin=*]
\item an explicit target predicate $P$ not already contained in $\mathcal{R}$;
\item a typed construction or estimator witnessing the objects on which $P$ is evaluated;
\item a proof, reduction, or reproducible protocol connecting $\mathcal{H}$ to $P$;
\item negative controls or failure modes under which $P$ must not be asserted;
\item a provenance statement distinguishing theorem, implementation, internal support, and external adjudication.
\end{enumerate}
\end{definition}

\begin{proposition}[No automatic downstream promotion]
Let $\mathfrak{E}=(\mathcal{O},\mathcal{H},\mathcal{R},\mathcal{B})$ be a BaseCore export object and let $P$ be a predicate introduced outside BaseCore. If $P\notin\mathcal{R}$, then BaseCore does not entail $P$ unless an admissible extension supplies an additional bridge from $\mathcal{H}$ to $P$.
\end{proposition}

\begin{proof}
By definition, $\mathcal{R}$ records the results proved from $\mathcal{H}$ inside BaseCore. A predicate $P$ introduced downstream is not in that result set. Therefore asserting $P$ from BaseCore alone would add an unstated inference. The missing inference must be supplied as a separate bridge, proof, reduction, or reproducible protocol.
\end{proof}

\begin{criterion}[Runtime-use admissibility]
A runtime measurement may instantiate a BaseCore export object only as an implementation-level witness when all of the following are recorded: the measured state space, the estimator or update rule, the finite horizon, the calibration regime, the controls under which the measurement fails, and the claim class being asserted. Such a measurement is not a BaseCore theorem and is not external validation by itself.
\end{criterion}

\begin{remark}[Open-load placement for $I_{\mathrm{int}}$ and atomic separators]
Quantities such as $I_{\mathrm{int}}$ and any atomic-separator condition may be treated as downstream target predicates or estimator families only after their carrier objects, separator class, invariance target, and negative controls are specified. Until a proof or reproducible adversarial protocol closes that burden, BaseCore may export the ambient structural language but not the separator conclusion.
\end{remark}

\begin{proposition}[Boundary preservation under finite implementations]
Suppose a finite implementation produces diagnostics for a downstream predicate $P$ over a bounded horizon. If the implementation does not prove that the diagnostics are invariant over the relevant mathematical class, then the diagnostics remain implementation evidence for $P$ and do not promote $P$ to a BaseCore result.
\end{proposition}

\begin{proof}
The implementation supplies finite observations or computations. Promotion to a BaseCore result requires a class-level implication from the stated hypotheses to the predicate. Without that implication, the observations can support, falsify, or calibrate a downstream claim, but they do not alter the theorem inventory of BaseCore.
\end{proof}
